\documentclass[11pt,letterpaper]{article}
\usepackage[utf8]{inputenc}
\usepackage{amsmath,amssymb,amsfonts}
\usepackage[margin=0.55in]{geometry}
\usepackage{enumitem}
\usepackage{chemfig}
\usepackage{mol2chemfig}
\usepackage{array}

\setchemfig{atom sep=1.8em, bond offset=1pt}

\begin{document}

% ==================== PAGE 1 ====================
\noindent


\vspace{0.2cm}
\noindent SCH4U \hfill Name: \underline{\hspace{6.5cm}} \\
\hfill Thursday, January 8, 2026

\vspace{0.2cm}
\begin{center}
\textbf{\large Unit 5: Organic Chemistry Half-Unit Test A}
\end{center}

\vspace{0.1cm}
\noindent
\begin{tabular}{|p{3.5cm}|p{3.5cm}|p{3.5cm}|p{3.5cm}|}
\hline
Total: \hfill /30 & = \hfill /15 & Comm: \hfill /10 & LL: \hfill \\
\hline
\end{tabular}

\vspace{0.4cm}
\noindent
1. \textbf{[6 marks]} Identify the type of organic molecule(s) based on the functional group(s) in the following:
\vspace{0.1cm}

\noindent
\begin{tabular}{|m{8.5cm}|m{8.2cm}|}
\hline
\hfil \textbf{Structure} \hfil & \hfil \textbf{Type of Organic Molecule(s) Based on Functional Group(s)} \hfil \\
\hline
\vspace{0.2cm}
\hfil \chemfig{-[:30]-[:-30](=[6]O)-[:30](-[:90])-[:-30](-[:-90]F)-[:30]-[:-30]} \hfil & \\
\hfil \small 5-fluoro-4-methyl-3-heptanone \hfil & \\[2.0cm]
\hline
\vspace{0.2cm}
\hfil \chemfig{=[:30]-[:-30]-[:30](-[:90]NH_2)=[:-30]-[:30]-[:-30]-[:30]} \hfil & \\
\hfil \small 4-amino-1,5-octadiene \hfil & \\[2.0cm]
\hline
\vspace{0.2cm}
\hfil \chemfig{-[:30](-[:90])-[:-30]((=[6]O))-[:30]O-[:-30]*6(=-=-=-)} \hfil & \\
\hfil \small phenyl \quad 2-methylpropanoate \hfil & \\[2.0cm]
\hline
\end{tabular}

\vspace{0.3cm}
\noindent
2. \textbf{[2 marks]} What effect does a double or triple bond have on the density of a compound? Explain.
\vspace{2.0cm}

\noindent
3. \textbf{[2 marks]} Which of the following structures should have the highest MP/BP? Explain.
\vspace{0.2cm}

\noindent
\begin{tabular}{|m{5.5cm}|m{5.5cm}|m{5.5cm}|}
\hline
\vspace{0.2cm} \hfil \chemfig{-[:30]=[:-30]-[:30]} \hfil & \vspace{0.2cm} \hfil \chemfig{-[:30]=[:-30]-[:-150]} \hfil & \vspace{0.2cm} \hfil \chemfig{*4(----)} \hfil \\ [0.2cm]
\hline
\hfil molecule X \hfil & \hfil molecule Y \hfil & \hfil molecule Z \hfil \\
\hline
\end{tabular}
\vspace{2.5cm}

\newpage
% ==================== PAGE 2 ====================
\noindent


\vspace{0.4cm}
\noindent
4. \textbf{What is the chemical formula of a structure that satisfies the following conditions? Show your work analytically as opposed to drawing a structure that fits the description.}

\begin{enumerate}[label=\alph*., leftmargin=0.5cm]
\item \textbf{[2 marks]} An aliphatic that has 40 carbon atoms, two cyclic structures \& has one triple bond.
\vspace{3.5cm}
\item \textbf{[2 marks]} An alcohol that has 20 carbon atoms, two hydroxyl groups \& two double bonds. No cyclic structure.
\end{enumerate}
\vspace{3.5cm}

\noindent
5. \textbf{[2 marks]} Would adding ethanamide \;\chemfig{-[:-30](=[6]O)-[:30]NH_2}\; to water cause boiling point elevation? Explain.
\vspace{3.5cm}

\noindent
6. \textbf{[2 marks]} Which of the following structures should have the best solubility in water? Explain.
\vspace{0.2cm}

\noindent
\begin{tabular}{|m{5.5cm}|m{5.5cm}|m{5.5cm}|}
\hline
\vspace{0.2cm} \hfil \textbf{CH}$_3$\textbf{OH} \hfil & \vspace{0.2cm} \hfil \chemfig{I-[:30]-[:-30]-[:30]-[:-30]-[:30]NH_2} \hfil & \vspace{0.2cm} \hfil \chemfig{-[:30](=[2]O)-[:-30]-[:30](=[2]O)-[:-30]O-[:30]-[:-30]} \hfil \\ [0.2cm]
\hline
\hfil molecule X \hfil & \hfil molecule Y \hfil & \hfil molecule Z \hfil \\
\hline
\end{tabular}
\vspace{3.5cm}

\newpage
% ==================== PAGE 3 ====================
\noindent


\vspace{0.4cm}
\noindent
7. \textbf{[3 marks]} Draw any 3 compounds that are structural isomers of the chemical formula $\text{C}_5\text{H}_8\text{O}$. Do not name the structures.
\vspace{4.0cm}

\noindent
8. \textbf{[2 marks]} Why do carboxylic acids tend to be solids or liquids at room temperature?
\vspace{2.0cm}

\noindent
9. \textbf{[2 marks]} Would methylcyclopropane \;\chemfig{-*3(---)}\; be stable at room temperature? Explain.
\vspace{2.5cm}

\noindent
10. \textbf{[2 marks]} Why do same mass aldehydes \& ketones have similar but slightly different MP/BP temperatures? Explain.
\vspace{2.5cm}

\noindent
11. \textbf{[1 mark]} Draw an example of a $3^\circ$ amine.
\vspace{2.0cm}

\noindent
12. \textbf{[2 marks]} Provide an explanation for the following BP data. Be specific.

\vspace{0.3cm}
\noindent
\begin{minipage}[t]{0.55\textwidth}
\vspace{0.1cm}
% Empty space for student responses
\vspace{3.0cm}
\end{minipage}
\hfill
\begin{minipage}[t]{0.4\textwidth}
\centering
\begin{tabular}{|c|c|}
\hline
\textbf{Molecule} & \textbf{BP Temperature ($^\circ$C)} \\
\hline
CH$_3$-F & -78.4 \\
\hline
CH$_3$-Cl & -24.2 \\
\hline
CH$_3$-Br & 3.56 \\
\hline
CH$_3$-I & 42.43 \\
\hline
\end{tabular}
\end{minipage}

\end{document}